\documentclass[11pt,a4paper]{article}

\usepackage[a4paper,margin=2.2cm]{geometry}
\usepackage[utf8]{inputenc}
\usepackage[T1]{fontenc}
\usepackage{booktabs}
\usepackage{xcolor}
\usepackage{hyperref}
\usepackage{amsmath}

\hypersetup{
    colorlinks=true,
    linkcolor=black,
    urlcolor=blue
}

\setlength{\parindent}{0pt}
\setlength{\parskip}{0.8em}
\setlength{\emergencystretch}{2em}

\begin{document}

\begin{center}
    {\Large Informationsvisualisierung -- Übung 7: Recursive-Pattern-Technik}
\end{center}

\section*{Kurzbeschreibung}

Die fünf Börsenindizes Dow Jones (DJ), Nikkei, Hangseng, DAX und Bovespa
werden mit der Recursive-Pattern-Technik visualisiert. Die Zeitreihen laufen
über verschiedene Zeiträume und werden einheitlich auf den Gesamtzeitraum
\textbf{1980-12-23 bis 2011-06-09} aufgefüllt. Die Implementierung umfasst
interaktive Indexauswahl sowie einen Umschalter zur Filterung der Wochenenden
(Aufgabe~7.7).

\section*{Lauffähiger Code}

\href{https://github.com/yaniktfl/iv-repo}{https://github.com/yaniktfl/iv-repo}

\section*{Ausführung}

Da Elm HTTP-Anfragen für lokale CSV-Dateien einen laufenden Webserver
benötigen, wird die Anwendung lokal betrieben:

\begin{verbatim}
elm make Main.elm --output=index.html
python3 -m http.server 9000
\end{verbatim}

Anschließend im Browser: \texttt{http://localhost:9000}

\section*{Aufgabe 7.1: CSV-Dateien laden}

Alle fünf CSV-Dateien werden beim Start der Anwendung parallel über HTTP
geladen. Dazu wird in \texttt{init} mit \texttt{Cmd.batch} eine Liste von
Commands erzeugt:

\begin{verbatim}
init _ =
    ( { selected = "DJ", weekendMode = WithWeekends
      , loaded = Dict.empty, failed = [] }
    , files |> List.map loadSeries |> Cmd.batch
    )
\end{verbatim}

Die Dateinamen werden als Liste von \texttt{(Name, Datei)}-Tupeln verwaltet.
Jedes Tupel erzeugt via \texttt{Http.get} ein eigenes Command. Die
\texttt{update}-Funktion empfängt die Antworten über die Message
\texttt{GotCsv String (Result Http.Error String)}.

\section*{Aufgabe 7.2: CSV-Parsing}

Aus jeder CSV-Datei werden die Felder \texttt{Date} und \texttt{Open}
extrahiert. Der Datentyp einer Tageszeile ist:

\begin{verbatim}
type alias StockDay = ( String, Maybe Float )
\end{verbatim}

Das Parsing nutzt \texttt{lovasoa/elm-csv} zum Tokenisieren und
\texttt{ericgj/elm-csv-decode} zum Dekodieren:

\begin{verbatim}
csvStringToData : String -> List StockDay
csvStringToData csvRaw =
    Csv.parse csvRaw
        |> Csv.Decode.decodeCsv decodeStockDay
        |> Result.toMaybe
        |> Maybe.withDefault []
\end{verbatim}

Fehlende oder nicht parsbare Kurswerte werden als \texttt{Nothing} gespeichert.

\section*{Aufgabe 7.3: Fehlende Tage auffüllen}

Da die Indizes unterschiedliche Startdaten haben und an Feiertagen keine Daten
liefern, werden alle Zeitreihen auf den gemeinsamen Zeitraum erweitert. Dazu
werden drei Schritte ausgeführt:

\begin{enumerate}
    \item Alle Kalendertage von 1980-12-23 bis 2011-06-09 werden mit
          \texttt{Date.range} erzeugt und als \texttt{Dict} mit
          \texttt{Nothing}-Werten initialisiert (\emph{Feiertagsgerüst}).
    \item Die tatsächlichen Börsendaten werden ebenfalls in ein \texttt{Dict}
          umgewandelt, gefiltert auf Tage die im Zielzeitraum liegen.
    \item \texttt{Dict.union actualDays emptyDays} liefert das Ergebnis:
          vorhandene Kurswerte bleiben erhalten, fehlende Tage erhalten
          \texttt{Nothing}.
\end{enumerate}

\begin{verbatim}
fillMissingDays : List String -> List StockDay -> List StockDay
fillMissingDays targetDays data =
    let
        emptyDays = targetDays
            |> List.map (\date -> ( date, Nothing ))
            |> Dict.fromList
        actualDays = data
            |> List.filter (\( date, _ ) -> Dict.member date emptyDays)
            |> Dict.fromList
    in
    Dict.union actualDays emptyDays |> Dict.toList
\end{verbatim}

\section*{Aufgaben 7.4 und 7.5: Recursive-Pattern zeichnen}

\subsection*{Pixelreihenfolge}

Anstelle des externen Pakets \texttt{elm-recursive-pattern} wurde
\texttt{createPixelMap} direkt implementiert. Die Funktion nimmt eine Liste
von Leveln und berechnet rekursiv die Pixel\-positionen:

\begin{verbatim}
type alias Level = { width : Int, height : Int }
type alias PixelPosition = { x : Int, y : Int }

createPixelMap : List Level -> List PixelPosition
\end{verbatim}

Jedes Level kachelt die inneren Positionen \texttt{(width $\times$ height)}-fach
nebeneinander. Die Level werden von außen nach innen definiert; die innersten
Level stehen am Ende der Liste.

\subsection*{Level-Konfiguration}

\begin{verbatim}
levels =
    [ Level 5 1,  Level 1 12
    , Level 4 1,  Level 2 3
    , Level 3 3,  Level 1 1  ]
\end{verbatim}

Das ergibt eine Gesamtgröße von
$5 \times 1 \times 4 \times 2 \times 3 \times 1 = 120$ Spalten und
$1 \times 12 \times 1 \times 3 \times 3 \times 1 = 108$ Zeilen,
also $12\,960$ Pixel -- genug für den gesamten Zeitraum mit Wochenenden.

\subsection*{Darstellung}

Die Pixelpositionen werden mit \texttt{List.map2} mit den Börsendaten
verknüpft. Jeder Datenpunkt wird als \texttt{rect} mit der berechneten
Position und einer aus dem Kurswert abgeleiteten Farbe gezeichnet:

\begin{verbatim}
drawCell : List Float -> RecordedData -> Svg Msg
\end{verbatim}

Die Farbskala ist \texttt{Scale.Color.viridisInterpolator} mit linearer
Normierung auf $[\min, \max]$ der vorhandenen Kurswerte. Fehlende Tage
erhalten ein neutrales Grau (\texttt{rgb 188 198 190}).

\section*{Aufgabe 7.6: Verschiedene Indizes}

Über eine Toolbar können alle fünf Indizes per Klick ausgewählt werden. Der
aktive Index wird schwarz hervorgehoben. Die Farbskala Viridis wurde gewählt,
weil sie perceptually uniform ist und auch für Farbenblinde gut unterscheidbar
bleibt.

\section*{Aufgabe 7.7: Wochenenden entfernen}

Ein Umschaltknopf wechselt zwischen Kalender- und Handelstag-Modus. Im
Handelstag-Modus werden Samstage und Sonntage aus der Zieldatumsreihe
gefiltert:

\begin{verbatim}
isTradingWeekday : Date.Date -> Bool
isTradingWeekday date =
    case Date.weekday date of
        Sat -> False
        Sun -> False
        _   -> True
\end{verbatim}

Dadurch reduziert sich die Pixelanzahl von $12\,960$ auf rund $9\,288$.
Die Level-Konfiguration wird entsprechend angepasst (\texttt{weekdayLevels}),
sodass die Gesamtgröße zur neuen Datenmenge passt.

\section*{Aufgabe 7.8: Level-Anpassung}

\begin{center}
\begin{tabular}{lll}
\toprule
Modus & Level-Liste & Bedeutung \\
\midrule
Mit Wochenenden    & $5\times1,\;1\times12,\;4\times1,\;2\times3,\;3\times3,\;1\times1$ & Monate, Wochen, Tage \\
Ohne Wochenenden   & $5\times1,\;1\times9,\;4\times1,\;2\times3,\;3\times3,\;1\times1$  & 5 Handelstage/Woche \\
\bottomrule
\end{tabular}
\end{center}

Im Handelstag-Modus gruppiert das zweite Level 9 statt 12 Wochen pro Block,
was die Aufteilung in Quartale erleichtert und die grauen Feiertags-Pixel
eliminiert.

\end{document}
